% ASSERT_CONVENTION: natural_units=natural, metric_signature=mostly_minus, fourier_convention=physics, coupling_convention=alpha_s, generator_normalization=T(fund)=1/2, covariant_derivative_sign=D_mu=partial_mu-igA
%
% Lecture 5: Tumbling, the MAC Criterion, and Complementarity
% Part of: The Dualised Standard Model -- Part III Lecture Notes
%
% Non-SUSY strong dynamics conventions (continued from Lectures 1-4):
%   T(fund) = 1/2 per Weyl fermion
%   b_0 = (11N_c - 2N_f)/3 for N_f Dirac flavours
%   C_2(fund) = (N^2-1)/(2N), C_2(adj) = N
%   C_2(sym_2) = (N+2)(N-1)/N, C_2(antisym_2) = (N-2)(N+1)/N
%   Delta C_2 = C_2(R_1) + C_2(R_2) - C_2(R_c) [MAC criterion, Raby-Dimopoulos-Susskind 1980]
%   A(fund) = 1, A(antifund) = -1 (cubic anomaly coefficient)
%
\documentclass[11pt,a4paper]{article}
% ASSERT_CONVENTION: natural_units=natural, metric_signature=mostly_minus, fourier_convention=physics, coupling_convention=alpha_s, generator_normalization=T(fund)=1/2, covariant_derivative_sign=D_mu=partial_mu-igA
%
% CONVENTIONS (Seiberg hep-th/9411149)
% Metric: (+,-,-,-) (mostly minus)
% T(fund) = 1/2 per Weyl fermion
% b_0 = (11N_c - 2N_f)/3 for N_f Dirac flavours
% D_mu = partial_mu - ig A_mu^a T^a
% A(fund) = 1 (cubic anomaly coefficient per fundamental Weyl fermion)
% alpha_s = g_s^2 / (4 pi)
% R[Q] = (N_f - N_c)/N_f; R[W] = 2
% N_f counts Dirac flavour pairs (each = 2 Weyl fermions)
%
% This file is \input{} by all lecture files.

%% ---------- Document class ----------
% (loaded by the wrapper; preamble only provides packages and macros)

%% ---------- Standard packages ----------
\usepackage{amsmath,amssymb,amsthm,mathtools}
\usepackage{hyperref}
\usepackage[capitalise,noabbrev]{cleveref}
\usepackage{enumitem}
\usepackage{booktabs}
\usepackage{array}
\usepackage{xcolor}
\usepackage{tikz}
\usetikzlibrary{arrows.meta,decorations.markings,positioning,calc}
\usepackage{fancybox}
\usepackage{tcolorbox}
\tcbuselibrary{skins,breakable}
\usepackage{slashed}              % Feynman slash notation
\usepackage{cite}                 % compressed citation lists

%% ---------- Hyperref configuration ----------
\hypersetup{
  colorlinks=true,
  linkcolor=blue!70!black,
  citecolor=green!50!black,
  urlcolor=blue!80!black,
  pdfauthor={Part III Lecture Notes},
  pdftitle={The Dualised Standard Model}
}

%% ---------- Theorem environments ----------
\theoremstyle{plain}
\newtheorem{theorem}{Theorem}[section]
\newtheorem{lemma}[theorem]{Lemma}
\newtheorem{proposition}[theorem]{Proposition}
\newtheorem{corollary}[theorem]{Corollary}

\theoremstyle{definition}
\newtheorem{definition}[theorem]{Definition}
\newtheorem{example}[theorem]{Example}
\newtheorem{exercise}{Exercise}[section]

\theoremstyle{remark}
\newtheorem{remark}[theorem]{Remark}

%% ---------- Physics macros: gauge groups and representations ----------
\newcommand{\Nc}{N_{\!c}}
\newcommand{\Nf}{N_{\!f}}
\newcommand{\SU}[1]{\mathrm{SU}(#1)}
\newcommand{\UU}[1]{\mathrm{U}(#1)}

% Representations
\newcommand{\fund}{\square}
\newcommand{\antifund}{\overline{\square}}
\newcommand{\adj}{\mathrm{adj}}
\newcommand{\antisym}{\wedge^{\!2}}        % antisymmetric rank-2
\newcommand{\sym}{\mathrm{S}_2}           % symmetric rank-2

%% ---------- Physics macros: group theory constants ----------
% Dynkin index: Tr(T^a T^b) = T(R) delta^{ab}
\newcommand{\Tfund}{\tfrac{1}{2}}          % T(fund) = 1/2 per Weyl fermion
\newcommand{\Tadj}{\Nc}                    % T(adj) = N_c
\newcommand{\Cfund}{\frac{\Nc^2 - 1}{2\Nc}}  % C_2(fund)
\newcommand{\Cadj}{\Nc}                    % C_2(adj)

% Cubic anomaly coefficient: A(R) = Tr_R[T^a {T^b, T^c}] with A(fund) = 1
\newcommand{\Afund}{1}

%% ---------- Physics macros: beta function ----------
\newcommand{\bzero}{b_0}
\newcommand{\bone}{b_1}
\newcommand{\alphas}{\alpha_s}
\newcommand{\gs}{g_s}

%% ---------- Physics macros: SUSY / R-charge ----------
\newcommand{\Rcharge}[1]{R[#1]}
\newcommand{\Wsup}{W}                      % superpotential

%% ---------- Physics macros: covariant derivative and field strength ----------
\newcommand{\Dmu}{D_\mu}
\newcommand{\Fmn}{F_{\mu\nu}}
\newcommand{\Ftilde}{\widetilde{F}_{\mu\nu}}

%% ---------- Physics macros: common abbreviations ----------
\newcommand{\ie}{\textit{i.e.}}
\newcommand{\eg}{\textit{e.g.}}
\newcommand{\cf}{\textit{cf.}}
\newcommand{\IR}{\ensuremath{\mathrm{IR}}}
\newcommand{\UV}{\ensuremath{\mathrm{UV}}}
\newcommand{\AF}{\ensuremath{\mathrm{AF}}}
\newcommand{\BZ}{\ensuremath{\mathrm{BZ}}}
\newcommand{\NSVZ}{\ensuremath{\mathrm{NSVZ}}}
\newcommand{\SQCD}{\ensuremath{\mathrm{SQCD}}}
\newcommand{\QCD}{\ensuremath{\mathrm{QCD}}}
\newcommand{\SM}{\ensuremath{\mathrm{SM}}}
\newcommand{\Tr}{\mathrm{Tr}}

%% ---------- Convention box macro ----------
\newtcolorbox{conventionbox}[1][]{%
  enhanced,
  colback=blue!5!white,
  colframe=blue!60!black,
  fonttitle=\bfseries,
  title={Conventions},
  #1
}

%% ---------- Comparison table style ----------
\newtcolorbox{comparisontable}[1][]{%
  enhanced,
  colback=orange!5!white,
  colframe=orange!70!black,
  fonttitle=\bfseries,
  title={Comparison},
  #1
}

%% ---------- Warning / important box ----------
\newtcolorbox{warningbox}[1][]{%
  enhanced,
  colback=red!5!white,
  colframe=red!70!black,
  fonttitle=\bfseries,
  title={Important},
  #1
}

%% ---------- Preview / forward-reference box ----------
\newtcolorbox{previewbox}[1][]{%
  enhanced,
  colback=green!5!white,
  colframe=green!50!black,
  fonttitle=\bfseries,
  title={Preview},
  #1
}

%% ---------- Page layout ----------
\usepackage[margin=2.5cm]{geometry}
\setlength{\parskip}{0.4em}
\setlength{\parindent}{0em}

%% ---------- Section formatting ----------
\usepackage{titlesec}
\titleformat{\section}{\Large\bfseries}{\thesection}{1em}{}
\titleformat{\subsection}{\large\bfseries}{\thesubsection}{1em}{}

%% ---------- End of preamble ----------


\title{\textbf{Lecture 5: Tumbling, the MAC Criterion, and Complementarity}}
\author{The Dualised Standard Model --- Part III Lecture Notes}
\date{Lent Term 2027}

\begin{document}
\maketitle

\tableofcontents
\newpage

%% ========================================================================
%%  CONVENTIONS
%% ========================================================================

\begin{conventionbox}[title={Conventions for Lecture 5}]
All conventions from Lectures~1--4 remain in force.  Additional
conventions specific to non-perturbative strong dynamics:

\renewcommand{\arraystretch}{1.3}
\begin{tabular}{@{}l l l@{}}
\toprule
\textbf{Quantity} & \textbf{Convention} & \textbf{Comment} \\
\midrule
Units & $\hbar = c = 1$ (natural) & \\
Metric & $\eta_{\mu\nu} = \mathrm{diag}(+1,-1,-1,-1)$ & mostly minus \\
Dynkin index & $\Tr(T^a T^b) = \Tfund\, \delta^{ab}$ & fundamental of $\SU{\Nc}$ \\
Cubic anomaly & $A(\fund) = 1$, $A(\antifund) = -1$ & per fundamental Weyl \\
Quadratic Casimir & $C_2(\fund) = \Cfund$ & fundamental \\
 & $C_2(\adj) = \Cadj$ & adjoint \\
 & $C_2(\antisym) = \dfrac{(\Nc-2)(\Nc+1)}{\Nc}$ & antisymmetric rank 2 \\[4pt]
 & $C_2(\sym) = \dfrac{(\Nc+2)(\Nc-1)}{\Nc}$ & symmetric rank 2 \\[4pt]
MAC attraction & $\Delta C_2 = C_2(R_1) + C_2(R_2) - C_2(R_c)$ &
  Raby--Dimopoulos--Susskind \\
\bottomrule
\end{tabular}
\end{conventionbox}


%% ========================================================================
%%  SECTION 1: Chiral Gauge Theories
%% ========================================================================
\section{Chiral Gauge Theories and Strong Dynamics}
\label{sec:chiral}

In Lectures~1--4 we studied \emph{vector-like} gauge theories: theories
in which the left-handed and right-handed fermions transform in
representations that are related by complex conjugation.  QCD is the
canonical example: each quark flavour comes with a Dirac partner, and
the gauge representation $\fund \oplus \antifund$ is real (in the sense
of being self-conjugate as a combined representation).  In such theories,
the standard expectation is chiral symmetry breaking through a
$\langle \bar{\psi}\psi \rangle$ condensate that preserves the gauge
symmetry.

In this lecture, we turn to a richer class of theories: \textbf{chiral
gauge theories}, in which the fermion content does \emph{not} form a
real representation of the gauge group.  The Georgi--Glashow $\SU{5}$
model with fermions in the $\mathbf{10} \oplus \overline{\mathbf{5}}$
representation is a classic example: the $\mathbf{10}$ (antisymmetric
rank-2) is not the conjugate of the $\overline{\mathbf{5}}$, and the
combined representation is irreducibly complex.

In chiral gauge theories, a $\langle \bar{\psi}\psi \rangle$ condensate
of the QCD type is not available---such a condensate would break the
gauge symmetry itself.  Instead, the strong dynamics can lead to
\textbf{dynamical gauge symmetry breaking} through condensation of
gauge-invariant fermion bilinears in specific channels.  Which channel
condenses first?  This is the question addressed by the Most Attractive
Channel (MAC) criterion.

\medskip

\textbf{Plan for this lecture.}  We develop three tools for analysing
strong dynamics in gauge theories:
\begin{enumerate}
  \item The \textbf{MAC criterion} (\S\ref{sec:MAC}): a heuristic for
    identifying which fermion bilinear condenses first.
  \item \textbf{Tumbling} (\S\ref{sec:tumbling}): the mechanism of
    sequential condensation and gauge group breaking, illustrated by
    the $\SU{5}$ Georgi--Glashow model.
  \item \textbf{Complementarity} (\S\ref{sec:complementarity}): the
    Fradkin--Shenker theorem relating the Higgs and confinement phases.
\end{enumerate}
In \S\ref{sec:connection} we explain how these tools connect to the
electric--magnetic duality framework from Lectures~3--4, setting up the
non-SUSY extensions in Lecture~6.


%% ========================================================================
%%  SECTION 2: The MAC Criterion
%% ========================================================================
\section{The Most Attractive Channel (MAC) Criterion}
\label{sec:MAC}

\subsection{The one-gluon exchange argument}

Consider a gauge theory with gauge group $\SU{\Nc}$ and fermions in
representations $R_1$ and $R_2$.  At one-loop order in perturbation
theory, the exchange of a single gluon between two fermions produces an
attractive potential in the colour channel $R_c$ that appears in the
tensor product $R_1 \otimes R_2$.  The strength of this attraction is
proportional to the \textbf{quadratic Casimir difference}:
\begin{equation}\label{eq:DeltaC2}
  \boxed{\Delta C_2 \;=\; C_2(R_1) + C_2(R_2) - C_2(R_c)\,.}
\end{equation}

The channel $R_c$ with the \textbf{largest} $\Delta C_2$ experiences the
strongest attraction.  The MAC hypothesis~\cite{RDS1980} asserts that
this most attractive channel is the one in which a fermion bilinear
condensate forms first as the gauge coupling runs to strong coupling.

\subsection{Physical picture}

The argument has a simple physical origin.  In the non-relativistic
limit, the one-gluon exchange potential between two colour sources in
representations $R_1$ and $R_2$, forming a composite in channel $R_c$,
is:
\begin{equation}\label{eq:Voge}
  V(r) \;\sim\; -\frac{\alphas}{r}\,
  \bigl[C_2(R_1) + C_2(R_2) - C_2(R_c)\bigr]
  \;=\; -\frac{\alphas}{r}\,\Delta C_2\,.
\end{equation}
A larger $\Delta C_2$ means a deeper attractive potential well, which
favours the formation of a bound state (condensate) in that channel.
The condensate with the strongest attraction forms first as the coupling
increases.

\begin{remark}[Why the smallest $C_2(R_c)$?]
  Since $\Delta C_2 = C_2(R_1) + C_2(R_2) - C_2(R_c)$ and $C_2(R_1)$,
  $C_2(R_2)$ are fixed by the fermion content, $\Delta C_2$ is maximised
  when $C_2(R_c)$ is \emph{minimised}.  In a tensor product
  $R_1 \otimes R_2 = \bigoplus_c R_c$, the smallest $C_2(R_c)$ typically
  corresponds to the smallest (lowest-dimensional) representation in the
  decomposition.  For $\fund \otimes \fund$, this is the singlet (if
  available) or the antisymmetric representation.
\end{remark}

\subsection{Quadratic Casimir values}

For reference, we collect the quadratic Casimirs for the common
representations of $\SU{N}$, using the conventions locked in Lecture~1
(with $T(\fund) = \tfrac{1}{2}$):

\begin{equation}\label{eq:casimirs}
\renewcommand{\arraystretch}{1.5}
\begin{array}{l c l}
\toprule
\textbf{Representation} & \textbf{Dimension} & C_2(R) \\
\midrule
\text{Fundamental } \fund & N & \dfrac{N^2 - 1}{2N} \\[6pt]
\text{Antifundamental } \antifund & N & \dfrac{N^2 - 1}{2N} \\[6pt]
\text{Adjoint } \adj & N^2 - 1 & N \\[6pt]
\text{Antisymmetric } \antisym & \dfrac{N(N-1)}{2} &
  \dfrac{(N-2)(N+1)}{N} \\[6pt]
\text{Symmetric } \sym & \dfrac{N(N+1)}{2} &
  \dfrac{(N+2)(N-1)}{N} \\
\bottomrule
\end{array}
\end{equation}

\textbf{Verification for $\SU{3}$:}
$C_2(\fund) = (9-1)/6 = 4/3$.
$C_2(\adj) = 3$.
$C_2(\antisym) = (1)(4)/3 = 4/3$.
(This is correct: for $\SU{3}$, the antisymmetric rank-2 representation
is the $\overline{\mathbf{3}}$, which is the antifundamental, so
$C_2(\antisym) = C_2(\fund)$.)\quad$\checkmark$


\subsection{Limitations of the MAC criterion}

\begin{warningbox}[title={MAC is a heuristic, not a theorem}]
The MAC criterion has significant limitations that must be stated
honestly:
\begin{enumerate}[label=(\alph*)]
  \item \textbf{MAC is a conjecture, not a theorem.}  There is no
    rigorous proof that the most attractive channel (as identified by
    one-gluon exchange) is the one that actually condenses in the
    full non-perturbative theory.

  \item \textbf{Perturbative reasoning at strong coupling.}  The
    one-gluon exchange argument~\eqref{eq:Voge} is a leading-order
    perturbative calculation.  However, condensation occurs in the
    strong-coupling regime ($\alphas \sim 1$), where higher-order
    corrections are not small and the perturbative picture is
    unreliable.

  \item \textbf{Known counterexamples.}  When two or more channels
    have comparable $\Delta C_2$ values, the MAC prediction can fail.
    Kosower~\cite{Kosower1983} demonstrated that for certain $\SU{N}$
    theories, the competition between channels is not resolved by the
    leading-order criterion alone.

  \item \textbf{Channel, not scale.}  Even when the MAC correctly
    identifies the condensation channel, it does not predict the
    condensation \emph{scale}.  The scale is set by the full
    non-perturbative dynamics, which the one-gluon exchange does not
    capture.
\end{enumerate}
\textbf{Usage:}  The MAC criterion \emph{suggests} the condensation
channel.  We do not say ``the condensate forms in the MAC.''  When the
dominant channel is clearly separated from competitors (as in our
$\SU{5}$ example below), the criterion is most reliable.  Within SUSY,
the exact results of Seiberg duality (\emph{cf.}\ Lecture~3) supersede
the MAC approximation.
\end{warningbox}


%% ========================================================================
%%  SECTION 3: Tumbling -- the SU(5) worked example
%% ========================================================================
\section{Tumbling: The $\SU{5}$ Worked Example}
\label{sec:tumbling}

We now demonstrate tumbling in the classic Georgi--Glashow $\SU{5}$
model~\cite{RDS1980}, with fermion content:
\begin{equation}\label{eq:SU5-content}
  \psi^{ab} \;\in\; \mathbf{10}\;(\antisym \text{ of } \SU{5})\,,
  \qquad
  \chi_a \;\in\; \overline{\mathbf{5}}\;(\antifund \text{ of } \SU{5})\,,
\end{equation}
where $a, b = 1, \ldots, 5$ are $\SU{5}$ fundamental indices, and
$\psi^{ab} = -\psi^{ba}$.  Both $\psi$ and $\chi$ are left-handed
Weyl fermions.

\begin{remark}[Anomaly cancellation]
  This fermion content is anomaly-free.  The cubic $\SU{5}^3$ anomaly is:
  $A(\mathbf{10}) + A(\overline{\mathbf{5}})
  = A(\antisym) + A(\antifund)
  = (5-4) + (-1) = 1 - 1 = 0$.\quad$\checkmark$
  (Recall from Lecture~1 that $A(\antisym) = N - 4$ and
  $A(\antifund) = -1$.)
\end{remark}

\subsection{Step 1: Enumerate condensation channels}

We consider bilinears formed from $\psi^{ab} \in \mathbf{10}$.  The
tensor product of two antisymmetric rank-2 representations of $\SU{5}$
decomposes as:
\begin{equation}\label{eq:10x10}
  \mathbf{10} \otimes \mathbf{10}
  \;=\;
  \overline{\mathbf{5}} \;\oplus\; \overline{\mathbf{45}}
  \;\oplus\; \mathbf{50}\,.
\end{equation}
The $\overline{\mathbf{5}}$ channel corresponds to the contraction
$\langle \psi^{ab}\, \psi^{cd}\rangle \sim \epsilon^{abcde}$, which
produces a composite transforming as $\overline{\mathbf{5}}$.

We also consider the mixed bilinear $\mathbf{10} \otimes
\overline{\mathbf{5}}$:
\begin{equation}\label{eq:10x5bar}
  \mathbf{10} \otimes \overline{\mathbf{5}}
  \;=\;
  \mathbf{5} \;\oplus\; \mathbf{45}\,.
\end{equation}

\subsection{Step 2: Compute $\Delta C_2$ for each channel}

Using the Casimir values from~\eqref{eq:casimirs} with $N = 5$:
\begin{align}
  C_2(\mathbf{10})
  &= C_2(\antisym\text{ of }\SU{5})
  = \frac{(5-2)(5+1)}{5}
  = \frac{18}{5} = 3.6\,,
  \label{eq:C2-10} \\
  C_2(\overline{\mathbf{5}})
  &= C_2(\fund\text{ of }\SU{5})
  = \frac{25-1}{10}
  = \frac{12}{5} = 2.4\,.
  \label{eq:C2-5bar}
\end{align}

\textbf{Channel comparison:}
\begin{equation}\label{eq:MAC-channels}
\renewcommand{\arraystretch}{1.4}
\begin{array}{l c c c}
\toprule
\textbf{Bilinear} & \textbf{Channel } R_c
  & \Delta C_2 = C_2(R_1) + C_2(R_2) - C_2(R_c)
  & \textbf{Value} \\
\midrule
\mathbf{10} \times \mathbf{10}
  & \overline{\mathbf{5}}
  & \tfrac{18}{5} + \tfrac{18}{5} - \tfrac{12}{5}
  & \dfrac{24}{5} = 4.8 \\[8pt]
\mathbf{10} \times \overline{\mathbf{5}}
  & \mathbf{5}
  & \tfrac{18}{5} + \tfrac{12}{5} - \tfrac{12}{5}
  & \dfrac{18}{5} = 3.6 \\[8pt]
\mathbf{10} \times \mathbf{10}
  & \overline{\mathbf{45}}
  & \tfrac{18}{5} + \tfrac{18}{5} - C_2(\overline{\mathbf{45}})
  & < 4.8 \\[4pt]
\bottomrule
\end{array}
\end{equation}
The $\overline{\mathbf{45}}$ channel has $C_2 > C_2(\overline{\mathbf{5}})$,
so its $\Delta C_2$ is smaller.  The most attractive channel is:
\begin{equation}\label{eq:MAC-result}
  \boxed{%
    \text{MAC:}\quad
    \mathbf{10} \times \mathbf{10}
    \;\to\; \overline{\mathbf{5}}\,,
    \qquad
    \Delta C_2 = \frac{24}{5}\,.
  }
\end{equation}

The MAC suggests that the $\SU{5}$ theory first forms the condensate:
\begin{equation}\label{eq:condensate}
  \langle \psi^{ab}\,\psi^{cd}\rangle
  \;\sim\; \epsilon^{abcde}\, v_e\,,
\end{equation}
where $v_e$ is a vector in the fundamental representation space of
$\SU{5}$ that specifies the direction of condensation.

\subsection{Step 3: Gauge symmetry breaking}

The condensate~\eqref{eq:condensate} transforms as a
$\overline{\mathbf{5}}$ under $\SU{5}$.  A non-zero expectation value
for this composite breaks the gauge symmetry.  Choosing (by gauge
rotation) $v_e = v\,\delta_{e5}$, the unbroken subgroup is:
\begin{equation}\label{eq:SU5-breaking}
  \SU{5} \;\xrightarrow{\;\langle\psi\psi\rangle \neq 0\;}
  \;\SU{4}\,.
\end{equation}
The condensate picks out the fifth direction, breaking $\SU{5}$ to the
$\SU{4}$ that rotates the remaining four indices $a = 1,2,3,4$.

\textbf{This is tumbling:}  the gauge group has broken to a smaller
group.  The residual $\SU{4}$ theory now contains the \emph{surviving}
massless fermions---those components of the original $\psi$ and $\chi$
that did not acquire masses from the condensate---and the process
repeats.  The cascade of sequential condensation steps is the
\textbf{tumbling mechanism}~\cite{RDS1980}.

\subsection{Fermion decomposition under $\SU{4}$}

Under $\SU{5} \to \SU{4}$, the representations decompose as:
\begin{align}
  \mathbf{10}\;(\psi^{ab}) &\;\to\;
    \mathbf{6}\;(\psi^{ij}) \;\oplus\;
    \mathbf{4}\;(\psi^{i5})\,,
    \label{eq:10-decomp} \\
  \overline{\mathbf{5}}\;(\chi_a) &\;\to\;
    \overline{\mathbf{4}}\;(\chi_i) \;\oplus\;
    \mathbf{1}\;(\chi_5)\,,
    \label{eq:5bar-decomp}
\end{align}
where $i,j = 1,\ldots,4$ are $\SU{4}$ fundamental indices.  The
$\mathbf{6}$ is the antisymmetric rank-2 of $\SU{4}$, and the
$\mathbf{4}$ comes from the components $\psi^{i5}$ with one index in
the broken direction.

The condensate $\langle \psi^{ab}\,\psi^{cd}\rangle
\sim \epsilon^{abcd5}\,v$ involves the components $\psi^{ij}$
(forming the $\SU{4}$-invariant combination
$\epsilon^{ijkl}\,\psi_{ij}\,\psi_{kl}$).  The condensation gives
mass to a linear combination of fermions, while the orthogonal
combinations remain massless.

The surviving massless fermions in the $\SU{4}$ theory are:
\begin{equation}\label{eq:SU4-content}
  \SU{4}:\qquad
  \psi^{i5} \in \mathbf{4}\,,\qquad
  \chi_i \in \overline{\mathbf{4}}\,,\qquad
  \chi_5 \in \mathbf{1}\,.
\end{equation}
(The $\mathbf{6}$ of $\SU{4}$ from $\psi^{ij}$ condenses and the
corresponding fermions acquire mass from the gap equation.)

\subsection{Step 4: Anomaly matching at the tumbling step}
\label{ssec:anomaly-tumbling}

We now verify that 't~Hooft anomaly matching---the exact constraint
from Lecture~1, \S3---is satisfied at the first tumbling step.  This
provides a non-trivial consistency check on the fermion identification
above.

The original $\SU{5}$ theory has a global $\UU{1}$ symmetry
(generalised baryon number) under which the fermions carry charges:
\begin{equation}\label{eq:U1-charges}
  B[\psi] = +1\,,\qquad B[\chi] = -3\,.
\end{equation}
(The relative charge is fixed by requiring $B[\langle\psi\psi\rangle]
\neq 0$ and by the anomaly-free condition for the gauge group
$\SU{5}^2\,\UU{1}_B$: $\;T(\mathbf{10})\cdot(+1) +
T(\overline{\mathbf{5}})\cdot(-3) = \tfrac{3}{2}\cdot 1 +
\tfrac{1}{2}\cdot(-3) = 0$.)

The relevant anomaly coefficient for $\UU{1}_B$ is the mixed
gauge--global anomaly $\SU{5}^2\,\UU{1}_B$, which vanishes by
construction.  However, we can check the \emph{purely global}
anomaly $\UU{1}_B^3$ as a non-trivial test.

\textbf{UV (before condensation):}
\begin{equation}\label{eq:U1B3-UV}
  \mathcal{A}[\UU{1}_B^3]_{\mathrm{UV}}
  = \dim(\mathbf{10})\cdot B[\psi]^3
  + \dim(\overline{\mathbf{5}})\cdot B[\chi]^3
  = 10 \cdot 1 + 5 \cdot (-27) = 10 - 135 = -125\,.
\end{equation}

\textbf{IR (after condensation):}  The surviving fermions
in~\eqref{eq:SU4-content} carry the same $\UU{1}_B$ charges as the
original fields from which they descend:
\begin{equation}\label{eq:U1B3-IR}
  \mathcal{A}[\UU{1}_B^3]_{\mathrm{IR}}
  = \dim(\mathbf{4})\cdot (1)^3
  + \dim(\overline{\mathbf{4}})\cdot (-3)^3
  + 1 \cdot (-3)^3
  = 4 - 108 - 27 = -131\,.
\end{equation}

The mismatch ($-125$ vs $-131$) indicates that this simple
decomposition does not yet account for all the massless fermion content
at the $\SU{4}$ stage---additional composite massless fermions
('t~Hooft's ``persistent mass condition'') must appear to satisfy
anomaly matching, or some of the fermion components from the condensed
representation remain massless as Goldstone partners.  This is
precisely the subtlety that makes tumbling analysis non-trivial:
anomaly matching \emph{constrains} the surviving spectrum but does not
uniquely determine it from the condensation channel alone.

\begin{remark}[The role of anomaly matching in tumbling]
  The point of this exercise is not that anomaly matching fails---it
  is an exact theorem and \emph{must} be satisfied---but that
  \emph{identifying the correct massless spectrum} at each tumbling
  step requires careful analysis beyond MAC alone.  The MAC tells us
  \emph{which channel condenses}; anomaly matching constrains
  \emph{which fermions remain massless}.  Together, they provide the
  tools for analysing the tumbling cascade.  We refer to the original
  analysis of Raby, Dimopoulos, and Susskind~\cite{RDS1980} and
  the pedagogical treatment in Peskin~\cite{Peskin1997} for the
  complete determination of the massless spectrum at each step of
  the $\SU{5}$ cascade.

  The key lesson from Lecture~1 applies here: \textbf{'t~Hooft anomaly
  matching does not determine the dynamics, but it powerfully constrains
  the possible low-energy spectra.}  Any proposed tumbling pattern that
  violates anomaly matching is ruled out.
\end{remark}


%% ========================================================================
%%  SECTION 4: Complementarity (Fradkin-Shenker)
%% ========================================================================
\section{Complementarity: The Fradkin--Shenker Theorem}
\label{sec:complementarity}

We now turn to a different but complementary perspective on the
relationship between the Higgs mechanism and confinement.

\subsection{Statement of the theorem}

Fradkin and Shenker~\cite{FS1979} (see also Osterwalder and
Seiler~\cite{OS1978}) proved a remarkable result on the lattice:

\begin{theorem}[Fradkin--Shenker, 1979]\label{thm:FS}
  For a gauge theory with a scalar field in the \textbf{fundamental
  representation} of the gauge group, the Higgs phase and the
  confinement phase are \textbf{analytically connected}:\ there is no
  phase transition separating them.
\end{theorem}

\textbf{Physical meaning.}  In a gauge theory with a fundamental
Higgs field $\phi$ (and no other order parameter distinguishing the
phases), one can continuously deform the theory from the
``Higgs regime'' (large scalar VEV, perturbative gauge boson masses)
to the ``confinement regime'' (strong coupling, flux tube formation)
without encountering a phase transition.  The two regimes are
\emph{not distinct phases}---they are different descriptions of the
same physics.

\subsection{The particle identification}

The physical consequence of complementarity is a \textbf{one-to-one
map} between the particle spectra in the two descriptions:

\renewcommand{\arraystretch}{1.4}
\begin{center}
\begin{tabular}{l l}
\toprule
\textbf{Higgs description} & \textbf{Confining description} \\
\midrule
Massive gauge bosons ($W$, $Z$) & Bound states of $\phi^\dagger\, D_\mu\, \phi$ \\
Higgs boson(s) & Bound states $\phi^\dagger \phi$ \\
Matter fermions $\psi$ & Confined composites $\phi^\dagger \psi$ \\
\bottomrule
\end{tabular}
\end{center}

The confined composites on the right are gauge-invariant operators
built from the scalar $\phi$ and the other fields.  Because there is
no phase transition, these gauge-invariant composites must match the
gauge-fixed particle descriptions on the left---their quantum numbers
(spin, flavour charges, baryon number, etc.)\ agree exactly.  This
identification is sometimes called the \textbf{complementarity
principle} for gauge theories~\cite{CT2011}.

\subsection{Conditions and limitations}

\begin{warningbox}[title={Complementarity is not universal}]
The Fradkin--Shenker theorem holds under two specific conditions:
\begin{enumerate}[label=(\alph*)]
  \item \textbf{The scalar must be in the fundamental representation}
    of the gauge group.

  \item \textbf{The gauge group must be completely broken} by the
    scalar VEV (no residual gauge symmetry).
\end{enumerate}
When either condition is violated, phase transitions \emph{can}
separate the Higgs and confinement phases:
\begin{itemize}
  \item For \textbf{adjoint Higgs} (or higher-representation scalars),
    the Higgs and confinement regimes can be separated by a genuine
    phase transition.  A discrete unbroken gauge symmetry (such as
    the centre $\mathbb{Z}_N$ of $\SU{N}$) provides a local order
    parameter that distinguishes the phases.

  \item For \textbf{partial gauge breaking}---where the scalar VEV
    leaves a residual gauge group---additional arguments are needed.
    The Standard Model electroweak sector has $\SU{2}_L \times
    \UU{1}_Y \to \UU{1}_{\mathrm{EM}}$, which is partial breaking.
    While complementarity-type arguments have been advanced for
    this case~\cite{CT2011}, the rigorous Fradkin--Shenker theorem
    does not directly apply.

  \item Shifman and Yung~\cite{SY2017} showed that in the
    Seiberg--Witten framework (Lecture~4), the quark confinement
    phase is \textbf{not} analytically connected to the Higgs phase.
    This provides a concrete counterexample to na\"{\i}ve universal
    complementarity when the conditions of the Fradkin--Shenker
    theorem are not met.
\end{itemize}
\end{warningbox}

\begin{remark}[Complementarity and the Standard Model]
  The SM Higgs doublet is in the fundamental ($\mathbf{2}$) of
  $\SU{2}_L$, satisfying condition~(a).  However, the electroweak
  breaking $\SU{2}_L \times \UU{1}_Y \to \UU{1}_{\mathrm{EM}}$
  is \emph{partial}: condition~(b) is not satisfied because
  $\UU{1}_{\mathrm{EM}}$ survives.  In the context of the dualised
  Standard Model (Lecture~6 and beyond), the Pati--Salam embedding
  provides a more natural setting where a larger gauge group is
  \emph{completely} broken, and complementarity applies more directly.
\end{remark}


%% ========================================================================
%%  SECTION 5: Connection to Duality
%% ========================================================================
\section{Connection to Duality}
\label{sec:connection}

\subsection{Tumbling + complementarity = a physical picture for duality}

We can now assemble the tools from this lecture into a \textbf{physical
picture} that motivates electric--magnetic duality:

\begin{enumerate}[label=\textbf{(\roman*)},leftmargin=*]
  \item \textbf{Electric theory (confining description).}  Start with
    a strongly-coupled gauge theory with fermions.  The MAC criterion
    suggests which fermion bilinear condenses first; the tumbling
    mechanism produces a cascade of condensation steps that break the
    gauge group sequentially.  At each step, massless composite
    fermions survive, constrained by anomaly matching.

  \item \textbf{Transition to the Higgs picture.}  At the end of the
    tumbling cascade, the strongly-coupled theory has produced
    condensates that break the gauge symmetry.  By complementarity,
    this confining description is continuously connected to a
    \emph{Higgs description} of the same physics: the condensates play
    the role of elementary Higgs fields, and the composites of the
    electric theory become the elementary particles of the Higgs phase.

  \item \textbf{Magnetic theory (Higgs description).}  The Higgs-phase
    description is a weakly-coupled theory with elementary scalars (the
    ``mesons'' of the electric theory) and gauge bosons---precisely the
    structure of the magnetic dual from Lecture~3.  The
    confined composites of the electric theory are identified with the
    elementary fields of the magnetic theory via their shared quantum
    numbers.
\end{enumerate}

\begin{remark}[From physical picture to duality]
  In the SUSY context of Lecture~3, this physical picture is made
  \emph{rigorous} by holomorphy, the NSVZ exact beta function, and
  anomaly matching.  In the non-SUSY context, we lose holomorphy and
  the exact beta function, but \textbf{anomaly matching survives}
  (it is a non-perturbative, UV/IR relation that does not require
  SUSY).  Tumbling and complementarity provide the physical intuition
  for \emph{why} a magnetic dual should exist; anomaly matching
  provides the \emph{constraint} that any proposed dual must satisfy.
\end{remark}

\subsection{What this lecture provides for Lecture~6}

The picture above---tumbling (dynamical mechanism) plus complementarity
(phase equivalence) producing a duality (theory equivalence)---is the
conceptual foundation for the non-SUSY duality programme.  In
Lecture~6 we will see how Sannino and collaborators used this framework,
combined with anomaly matching, to conjecture that the Standard Model
itself is the magnetic dual of a more fundamental electric theory.

\begin{previewbox}[title={Preview: Lecture 6}]
  In Lecture~6 we cross from established non-perturbative phenomena
  (tumbling, complementarity) to \textbf{conjectural} non-SUSY
  duality extensions.  The key question: can the physical picture of
  tumbling + complementarity, which is made rigorous by SUSY in
  Lectures~3--4, survive in the absence of supersymmetry?

  The evidence is compelling but not conclusive: anomaly matching
  works, perturbative fixed points exist, and the framework produces
  the correct Standard Model structure.  But non-SUSY duality has
  no independent cross-validation beyond anomaly matching---in
  contrast to the SUSY case, where holomorphy, moduli space matching,
  and the superconformal index all provide additional evidence.

  \textbf{Every non-SUSY duality statement in Lecture~6 will carry
  an explicit conjectural qualifier.}
\end{previewbox}


%% ========================================================================
%%  COMPARISON TABLE
%% ========================================================================
\section{Comparison: Tumbling, Complementarity, and Duality}
\label{sec:comparison}

\begin{comparisontable}[title={Three distinct concepts}]
These three ideas are often discussed together, but they are
\textbf{logically distinct}.  Conflating them is a common source of
confusion.

\renewcommand{\arraystretch}{1.5}
\begin{tabular}{@{} >{\raggedright}p{3cm} p{3.5cm} p{3.5cm} p{3.5cm} @{}}
\toprule
& \textbf{Tumbling} & \textbf{Complementarity} & \textbf{Duality} \\
\midrule
\textbf{What it is}
  & Dynamical mechanism
  & Phase-structure statement
  & Equivalence between theories \\
\textbf{Operates on}
  & A single theory
  & A single theory (two phases)
  & Two different theories \\
\textbf{Key idea}
  & Sequential condensation driven by MAC
  & Higgs phase = confinement phase (no phase transition)
  & Electric theory and magnetic theory have the same IR physics \\
\textbf{Evidence level}
  & Approximate (MAC heuristic)
  & Proven (for fundamental Higgs with complete breaking;
    \cite{FS1979})
  & Proven (SUSY, Lectures~3--4); conjectured (non-SUSY) \\
\textbf{Uses}
  & Identifies \emph{which} condensate forms
  & Identifies confined composites with Higgs-phase particles
  & Maps the full spectrum and interactions between two theories \\
\textbf{Limitations}
  & One-gluon exchange at strong coupling; counterexamples exist
  & Requires fundamental Higgs and complete breaking
  & Non-SUSY case unproven \\
\bottomrule
\end{tabular}
\end{comparisontable}


%% ========================================================================
%%  SUMMARY
%% ========================================================================
\section{Summary}
\label{sec:summary}

\begin{tcolorbox}[
  enhanced,
  colback=green!5!white,
  colframe=green!60!black,
  fonttitle=\bfseries,
  title={Lecture 5: Summary}
]
\begin{enumerate}
  \item \textbf{Chiral gauge theories} have fermion content that does
    not form a real representation; their strong dynamics can produce
    dynamical gauge symmetry breaking rather than the chiral symmetry
    breaking familiar from QCD.

  \item The \textbf{MAC criterion} identifies the most attractive
    condensation channel via
    $\Delta C_2 = C_2(R_1) + C_2(R_2) - C_2(R_c)$.
    It is a \emph{heuristic} based on one-gluon exchange, not a
    theorem.

  \item \textbf{Tumbling} is the sequential condensation mechanism:
    the MAC condenses, the gauge group breaks, and the process repeats.
    The $\SU{5}$ Georgi--Glashow model demonstrates the first step:
    $\SU{5} \to \SU{4}$ with $\Delta C_2 = 24/5$.

  \item At each tumbling step, \textbf{'t~Hooft anomaly matching}
    (from Lecture~1) constrains the surviving massless spectrum.
    Anomaly matching does not determine the dynamics, but any
    proposed tumbling pattern that violates it is excluded.

  \item \textbf{Fradkin--Shenker complementarity:}  for gauge theories
    with a fundamental-representation Higgs field and complete gauge
    breaking, the Higgs and confinement phases are analytically
    connected---there is no phase transition.  This is a
    \emph{theorem}, not a conjecture.

  \item \textbf{Limitations of complementarity:}  it does not apply
    universally.  For adjoint Higgs, partial gauge breaking, or
    higher-representation scalars, phase transitions can separate
    the phases.

  \item \textbf{Tumbling + complementarity} provides a physical
    picture motivating electric--magnetic duality: the confined
    composites of the electric theory are identified with the
    elementary fields of the magnetic (Higgs-phase) theory.
\end{enumerate}
\end{tcolorbox}


%% ========================================================================
%%  BIBLIOGRAPHY
%% ========================================================================
\begin{thebibliography}{99}

\bibitem{RDS1980}
  S.~Raby, S.~Dimopoulos, and L.~Susskind,
  ``Tumbling gauge theories,''
  \textit{Nucl.~Phys.} \textbf{B169} (1980) 373.

\bibitem{FS1979}
  E.~Fradkin and S.~H.~Shenker,
  ``Phase diagrams of lattice gauge theories with Higgs fields,''
  \textit{Phys.~Rev.} \textbf{D19} (1979) 3682.

\bibitem{OS1978}
  K.~Osterwalder and E.~Seiler,
  ``Gauge field theories on a lattice,''
  \textit{Ann.~Phys.} \textbf{110} (1978) 440.

\bibitem{Kosower1983}
  D.~A.~Kosower,
  ``Symmetry breaking patterns in pseudoreal and real gauge theories,''
  \textit{Phys.~Lett.} \textbf{B130} (1983) 141.

\bibitem{CT2011}
  C.~Cs\'aki and T.~Terning,
  ``Seiberg duality and the Standard Model,''
  [\texttt{arXiv:1106.3074}].

\bibitem{SY2017}
  M.~Shifman and A.~Yung,
  ``Fradkin--Shenker continuity and `instead-of-confinement' phase,''
  [\texttt{arXiv:1704.06201}].

\bibitem{Peskin1997}
  M.~E.~Peskin,
  ``Duality in supersymmetric Yang--Mills theory,''
  in \textit{Fields, Strings and Duality (TASI 96)},
  [\texttt{hep-th/9702094}].

\bibitem{Seiberg1995}
  N.~Seiberg,
  ``Electric--magnetic duality in supersymmetric non-Abelian gauge
  theories,''
  \textit{Nucl.~Phys.} \textbf{B435} (1995) 129
  [\texttt{hep-th/9411149}].

\bibitem{Sannino2024}
  F.~Sannino \textit{et al.},
  ``The Dualised Standard Model,''
  [\texttt{arXiv:2407.17281}].

\end{thebibliography}

\end{document}
